\documentclass[12pt]{article}
\usepackage{amsmath}
\usepackage{graphicx}
\usepackage{cite}
\usepackage[margin=1in]{geometry}
\usepackage{url}
\usepackage{hyperref}

\title{Study of GPU Kernel Optimizations for Scaled Dot-Product Attention}
\author{Piyush Jadhav}
\date{}

\begin{document}

\maketitle

\section{Introduction}

Transformer-based architectures have become the foundation of modern deep learning systems, particularly large language models (LLMs) used in natural language processing and generative AI \cite{vaswani2017attention}. At the core of these architectures lies scaled dot-product attention, an operation whose computational complexity grows quadratically with sequence length. As context windows and model sizes increase, attention becomes a major performance bottleneck during inference.

Recent research has shown that attention performance on GPUs is frequently constrained by memory bandwidth and data movement rather than peak floating-point throughput \cite{dao2022flashattention, wang2019characterizing}. IO-aware algorithmic restructuring techniques, such as FlashAttention \cite{dao2022flashattention}, significantly reduce high-bandwidth memory (HBM) transfers and improve efficiency. However, much of the literature proposes new algorithms or large-scale system optimizations, rather than isolating individual kernel-level techniques under controlled experimental settings.

In practice, GPU performance is highly sensitive to implementation details such as memory access patterns, shared memory utilization, kernel launch overhead, and occupancy. While techniques like shared-memory tiling and kernel fusion are commonly applied in high-performance GPU programming, there is limited empirical analysis examining when these optimizations meaningfully improve scaled dot-product attention performance and what tradeoffs they introduce.

This work presents a controlled performance analysis of two GPU kernel optimization techniques—shared-memory tiling and kernel fusion—applied to scaled dot-product attention during Transformer inference. By systematically varying workload configurations and analyzing hardware performance metrics, this study aims to provide insight into when these techniques are beneficial and when their impact diminishes.

\section{Background}

Scaled dot-product attention computes:

\[
\text{Attention}(Q,K,V) =
\text{softmax}\left(\frac{QK^T}{\sqrt{d}}\right)V
\]

where $Q$, $K$, and $V$ are query, key, and value tensors of shape $B \times N \times d$, with batch size $B$, sequence length $N$, and head dimension $d$ \cite{vaswani2017attention}. The computation involves two matrix multiplications and a row-wise softmax operation.

On GPUs, the $QK^T$ multiplication is typically implemented as a batched matrix multiplication, followed by scaling and normalization. The performance of these operations depends heavily on memory hierarchy usage, arithmetic intensity, and thread-level parallelism.

\section{Related Work}

\subsection{Transformer Attention}

The Transformer architecture introduced by Vaswani et al. \cite{vaswani2017attention} replaced recurrent neural networks with self-attention mechanisms. While highly parallelizable, scaled dot-product attention incurs $O(N^2)$ complexity in sequence length, motivating optimization efforts.

\subsection{IO-Aware Attention Optimization}

Dao et al. \cite{dao2022flashattention} demonstrated that attention performance is often memory-bound rather than compute-bound. FlashAttention reduces memory traffic by restructuring computation to maximize on-chip data reuse. FlashAttention-2 further improves parallelism and workload partitioning \cite{dao2023flashattention2}. These works emphasize that memory access patterns are critical to performance.

\subsection{GPU Bottleneck Characterization}

Wang et al. \cite{wang2019characterizing} performed architectural analysis of deep learning workloads on modern GPUs, showing that many operations operate near memory bandwidth limits. Their findings support investigating memory reduction techniques such as shared-memory tiling.

The Roofline model \cite{williams2009roofline} provides a theoretical framework for analyzing whether workloads are compute-bound or memory-bound. By relating arithmetic intensity to hardware ceilings, it enables systematic interpretation of performance measurements.

\subsection{ML Systems Optimization}

Jia et al. \cite{jia2019optimizing} emphasize systematic performance analysis of deep learning computation graphs and the importance of studying optimization tradeoffs rather than isolated improvements. This motivates the controlled experimental design adopted in this work.

\subsection{Gap in Existing Literature}

While prior research establishes that attention is memory-intensive and proposes algorithmic improvements, there is limited controlled empirical evaluation isolating kernel-level techniques such as shared-memory tiling and kernel fusion within standard scaled dot-product attention. This study addresses that gap by analyzing these techniques under varying workload configurations during inference.

\section{Methodology}

\subsection{Attention Formulation}

We study single-head scaled dot-product attention during inference,
defined as:

\begin{equation}
    \text{Attention}(Q, K, V) =
    \text{softmax}\!\left(\frac{QK^{\top}}{\sqrt{d}}\right)V
    \label{eq:attention}
\end{equation}

\noindent where $Q, K, V \in \mathbb{R}^{B \times N \times d}$,
$B$ is the batch size, $N$ is the sequence length, and $d$ is the
head dimension. The intermediate score matrix
$S = QK^{\top}/\sqrt{d} \in \mathbb{R}^{B \times N \times N}$
grows quadratically with sequence length, making memory bandwidth
the primary performance bottleneck at large $N$.

\subsection{Baseline Implementation}

The baseline kernel decomposes attention into three sequential
CUDA kernels, all operating exclusively on global memory:

\begin{enumerate}
    \item \textbf{QK$^{\top}$ with scaling} (\texttt{qk\_dot\_scaled\_kernel}):
    Computes $S_{b,i,j} = \sum_{k} Q_{b,i,k} \cdot K_{b,j,k} / \sqrt{d}$.
    Each thread computes one element of $S$. Grid: $(\lceil N/16 \rceil,
    \lceil N/16 \rceil, B)$, block: $(16, 16)$.

    \item \textbf{Row-wise softmax} (\texttt{softmax\_kernel}):
    Applies numerically stable softmax over the $j$-dimension of $S$
    using the max-subtraction trick. Each thread processes one complete
    row $S_{b,i,:}$. Grid: $(\lceil N/256 \rceil, B)$, block: $(256)$.

    \item \textbf{Weighted sum with V} (\texttt{attn\_output\_kernel}):
    Computes $\text{out}_{b,i,k} = \sum_{j} S_{b,i,j} \cdot V_{b,j,k}$.
    Each thread computes one output element. Grid:
    $(\lceil d/16 \rceil, \lceil N/16 \rceil, B)$, block: $(16, 16)$.
\end{enumerate}

This implementation is intentionally naive — no shared memory tiling,
no kernel fusion — to establish a clean performance baseline against
which optimized variants are compared in Sections~\ref{sec:tiling}
and~\ref{sec:fusion}.

\subsection{Correctness Validation}

Correctness was verified by comparing CUDA output against a PyTorch
reference implementation using identical input tensors. Input tensors
$Q$, $K$, $V$ were generated with uniform random values in $[-1, 1]$
using fixed seeds (42, 43, 44 respectively) to ensure reproducibility.
The CUDA binary writes output tensors to raw float32 binary files,
which are loaded in a validation notebook and compared element-wise
against the PyTorch reference (\texttt{torch.bmm} + \texttt{F.softmax}).

For the configuration $B{=}1$, $N{=}512$, $d{=}64$:

\begin{itemize}
    \item Maximum absolute error: $1.34 \times 10^{-7}$
    \item Mean absolute error: $1.13 \times 10^{-8}$
    \item Maximum relative error: $1.51 \times 10^{-6}$
\end{itemize}

These values are at the float32 machine epsilon floor
($\approx 1.19 \times 10^{-7}$), confirming numerical correctness
with no algorithmic errors. The tolerance threshold used is
$\epsilon = 10^{-3}$; all errors fall well below this bound.


\section{Experimental Setup}
\label{sec:setup}

\subsection{Hardware and Software Environment}

Experiments are conducted on a GPU cluster node equipped with an
NVIDIA \texttt{[TODO: GPU model, e.g. A100 / V100 / RTX 3090]}.
The CUDA toolkit version is \texttt{[TODO: nvcc --version]},
compiled with \texttt{-arch=sm\_[TODO]} and \texttt{--use\_fast\_math}.
The host system runs \texttt{[TODO: OS, e.g. Ubuntu 20.04]} with
\texttt{[TODO: CPU model]}. PyTorch \texttt{[TODO: version]} is used
solely for correctness validation and is not involved in any
performance measurements.

\subsection{Workload Configurations}

We evaluate all combinations of the following parameters:

\begin{table}[h]
\centering
\begin{tabular}{ll}
\hline
\textbf{Parameter} & \textbf{Values} \\
\hline
Batch size $B$       & 1, 4, 16, 32 \\
Sequence length $N$  & 128, 256, 512, 1024 \\
Head dimension $d$   & 64, 128 \\
\hline
\textbf{Total}       & \textbf{32 configurations} \\
\hline
\end{tabular}
\caption{Parameter sweep for baseline performance characterization.}
\label{tab:configs}
\end{table}

Input tensors $Q$, $K$, $V$ are generated synthetically using
uniform random values with fixed seeds, simulating real attention
workloads without requiring a dataset.

\subsection{Measurement Protocol}

Each configuration is executed with 5 warmup iterations followed by
100 timed iterations. Timing is performed using \texttt{cudaEvent}
start and stop records around the measured iterations, with
\texttt{cudaEventSynchronize} ensuring all kernels complete before
the elapsed time is read. The reported latency is the mean over 100
iterations. Throughput is derived as:

\begin{equation}
    \text{Throughput (tokens/sec)} = \frac{B \times N}{\text{latency (s)}}
\end{equation}

All results are written to a CSV file
(\texttt{results/week2/baseline\_sweep.csv}) with columns:
\texttt{variant}, \texttt{B}, \texttt{N}, \texttt{d},
\texttt{latency\_ms}, \texttt{throughput\_k\_tokens\_per\_sec}.
The same harness will be reused for tiled and fused variants in
subsequent phases, with the \texttt{variant} field distinguishing
implementations.